\section{Normalización física y estimación hasta grado logarítmico}
Fijados el exterior $A_0=A\setminus U$ y la cara $Y$, sea $\mathcal I_Y$ la familia de
subconjuntos de $U$ compatibles con RC y con las restricciones diagonales de $Y$.
Para $I\in\mathcal I_Y$, $m=j-|I|$. El átomo pesa $W_I=\operatorname{per}(A_{0,-I})$
en la máscara binaria. Se conservan los factores exactos de escalado si se usan matrices
balanceadas. Los menores utilizados tienen el escalado de Menores uniformes 02; un cero
impuesto se conserva como cero. Cada configuración contribuye una sola vez a la suma.

\paragraph{Los mismos pesos al tamaño físico.}
Escribamos $r=\nabla R(0)$, $R_c=R-R(0)-r\cdot x$. El peso que elegimos para los dos
aproximantes, en $s=1/j$, es
\begin{align*}
 \log(\widehat W_I/W_\varnothing)
 &=r\cdot{\bf1}_I+R_c({\bf1}_I)
   +A_*({\bf1}_I)-A_*(0)\\
 &\quad+\widetilde d_T(I)-\widetilde d_T(\varnothing),\qquad
 W_\varnothing=\operatorname{per}(A_0).
\end{align*}
$\widetilde d_T$ usa el mismo $q=4v+12$ en las dos evaluaciones, con el escalado y tamaño
propios de cada menor. En el átomo el margen apagado vale cero. Se llega por
$x=\lambda{\bf1}_I$ hasta $\lambda=1$ mediante las medidas levantadas. Por Menores
uniformes 07 los grados normalizados de los índices que desaparecen tienden a uno y el
cofactor sin covolumen gana $\tfrac12\log(m/j)$. Sumando $\tfrac12\log(4\pi j)$ se
obtiene exactamente $\tfrac12\log(4\pi m)$. La acción tiene el valor de torres del menor
porque los márgenes supervivientes valen uno y $t\log t$ vale cero en $t=0,1$.
Así se conserva el normalizador individual de cada $I$. La única constante al pasar del
logaritmo del gas al de la cara es $\log(W_\varnothing/\operatorname{per}A)$.
Ni $m$ ni un resto dependiente de $I$ se sacan de la suma.

A tamaño físico, Menores uniformes 01 y 05 dan
$|\log(\widehat W_I/W_I)|\le(Cq/j)^q+Ce^{-cj}$, aumentando $C$ para dos evaluaciones.
La versión con Stirling truncado y la de factorial exacto pueden diferir por átomo:
el mismo resto cubre esa diferencia. Márgenes y átomos 03 permite elegir el peso escrito
aquí. Si se fija la otra receta, su sustitución en la suma de diferencias cuesta
$O(j^{-8})$, por el lema de esa parte aplicado dos veces. Se aplica sólo en el tamaño
físico. En las caras las sumas reales son positivas. Su extensión holomorfa local procede
del gas construido para el aproximante, nunca de un Taylor en $s$ del error real.

\paragraph{De los coeficientes locales a la media.}
Como $\zeta_i=0$ implica $x_i=0$, las dos caras que difieren en la activación de $i$
tienen el mismo cofactor (Menores uniformes 07) y la misma suma de átomos.
Para $c$ componentes y rango cíclico $\beta$ en la incidencia completa ponemos
\[
 p=v+c+1\quad(\beta=0),\qquad p=v+c\quad(\beta=1),\qquad
 p=v+c-1\quad(\beta\ge2).
\]
Bajo los lemas previos de bosque, raíz y ciclo, la parte 06 de Menores uniformes puede
aplicarse ahora al exterior de las partes 01--02. Su conclusión a campos fijos y la
coincidencia de caras dan lo siguiente. En bosque el primer coste es $c+1$. Con un solo
ciclo, el coste $c-1$ conserva al menos un campo repetido y el resto empieza en $c$.
Con varios ciclos basta conexidad y el coste es $c-1$. El factor $2^v$ de las caras se
conserva. Cauchy en $s$ cuesta $(Cv/j)^w$ a partir del coste $w$, y en campos cuesta
$\prod_i(Cd_i/j)$; para la repetición se añade $\sum_i Cd_i/j\le Cv^2/j$.
El logaritmo $O(v)$ se absorbe en $C^v$. En las tres clases se obtiene
\[
 |\Delta_U(\widehat F_{U,q}-\Lambda)|
 \le C^v(Cv)^{c+1}\prod_{i\in U}d_i\, j^{-p}.
\]
Las constantes dependen de la clase de referencia, fijada antes de $b$ y del radio
global $\rho$. Aquí se aplican los lemas anteriores sobre el exterior común; la serie del
hipergrafo uniforme no interviene. Se excluyen las líneas individuales del resto entre
líneas. Para $v\ge5$ las sustracciones de parejas desaparecen.

Sea $H_j$ la llegada regular corta histórica, y $N=\lceil4\log n/\log\rho\rceil$.
Con la cuenta ponderada de tipos y las prescripciones de Llegada y completabilidad 03,
la contribución en grado $v$ está acotada por
\[
 (C_L/n)(C_1\delta^{-2})^v v^d\,\delta_*^{(v-4)/3},\qquad 5\le v\le N .
\]
La razón exponencial se fija menor que uno antes de tomar $b\le\delta_*$.
La suma de los valores absolutos en Márgenes y átomos 04 es $O_{b,\delta}(1/n)$
bajo la ley $\nu$ con indicador $H_j$. La comparación con los átomos exactos cuesta
$O(j^{-8})$ en el conjunto de todos esos grados. La implicación invoca las entradas
uniformes del mapa; no usa independencia de los índices en la historia.
